\documentclass[10pt,a4paper]{amsart}
\usepackage{amsmath,amssymb,amsthm,graphicx}
% Retained minimal preamble for the Stacks Project chapter more-algebra.tex (More on Algebra).
% The chapter file pinned-original-source.tex carries only \input{preamble} on its first line, so no
% preamble text lives inside the chapter itself.  The Stacks Project uses one single shared preamble for
% every chapter; the snapshot stages it identically as sources/stacks-chow/preamble.tex and
% sources/stacks-etale-cohomology/preamble.tex.  The block below is copied mechanically, line for line,
% from sources/stacks-chow/preamble.tex lines 14-18 (the Xy-pic setup) and 154-217 (the native theorem
% declarations together with the mathematical macros).  Nothing was added, removed or redefined; the
% wrapper class and the amsmath/amssymb/amsthm/graphicx base packages are supplied by the assembler.
% Chapter: more-algebra.tex -- More on Algebra.
% For commutative diagrams we use Xy-pic
\usepackage[all]{xy}

% We use 2cell for 2-commutative diagrams.
\xyoption{2cell}

% Theorem environments.
%
\theoremstyle{plain}
\newtheorem{theorem}[subsection]{Theorem}
\newtheorem{proposition}[subsection]{Proposition}
\newtheorem{lemma}[subsection]{Lemma}

\theoremstyle{definition}
\newtheorem{definition}[subsection]{Definition}
\newtheorem{example}[subsection]{Example}
\newtheorem{exercise}[subsection]{Exercise}
\newtheorem{situation}[subsection]{Situation}

\theoremstyle{remark}
\newtheorem{remark}[subsection]{Remark}
\newtheorem{remarks}[subsection]{Remarks}

\numberwithin{equation}{subsection}

% Macros
%
\def\lim{\mathop{\mathrm{lim}}\nolimits}
\def\colim{\mathop{\mathrm{colim}}\nolimits}
\def\Spec{\mathop{\mathrm{Spec}}}
\def\Hom{\mathop{\mathrm{Hom}}\nolimits}
\def\Ext{\mathop{\mathrm{Ext}}\nolimits}
\def\SheafHom{\mathop{\mathcal{H}\!\mathit{om}}\nolimits}
\def\SheafExt{\mathop{\mathcal{E}\!\mathit{xt}}\nolimits}
\def\Sch{\mathit{Sch}}
\def\Mor{\mathop{\mathrm{Mor}}\nolimits}
\def\Ob{\mathop{\mathrm{Ob}}\nolimits}
\def\Sh{\mathop{\mathit{Sh}}\nolimits}
\def\NL{\mathop{N\!L}\nolimits}
\def\CH{\mathop{\mathrm{CH}}\nolimits}
\def\proetale{{pro\text{-}\acute{e}tale}}
\def\etale{{\acute{e}tale}}
\def\QCoh{\mathit{QCoh}}
\def\Ker{\mathop{\mathrm{Ker}}}
\def\Im{\mathop{\mathrm{Im}}}
\def\Coker{\mathop{\mathrm{Coker}}}
\def\Coim{\mathop{\mathrm{Coim}}}

% Boxtimes
%
\DeclareMathSymbol{\boxtimes}{\mathbin}{AMSa}{"02}

%
% Macros for moduli stacks/spaces
%
\def\QCohstack{\mathcal{QC}\!\mathit{oh}}
\def\Cohstack{\mathcal{C}\!\mathit{oh}}
\def\Spacesstack{\mathcal{S}\!\mathit{paces}}
\def\Quotfunctor{\mathrm{Quot}}
\def\Hilbfunctor{\mathrm{Hilb}}
\def\Curvesstack{\mathcal{C}\!\mathit{urves}}
\def\Polarizedstack{\mathcal{P}\!\mathit{olarized}}
\def\Complexesstack{\mathcal{C}\!\mathit{omplexes}}
% \Pic is the operator that assigns to X its picard group, usage \Pic(X)
% \Picardstack_{X/B} denotes the Picard stack of X over B
% \Picardfunctor_{X/B} denotes the Picard functor of X over B
\def\Pic{\mathop{\mathrm{Pic}}\nolimits}
\def\Picardstack{\mathcal{P}\!\mathit{ic}}
\def\Picardfunctor{\mathrm{Pic}}
\def\Deformationcategory{\mathcal{D}\!\mathit{ef}}
\usepackage{cleveref}
\crefname{originallemma}{original lemma}{original lemmas}
\Crefname{originallemma}{Original lemma}{Original lemmas}
\makeatletter
\newlabel{more-algebra-lemma-finite-cohomology}{{064S}{source}{}{sourcepointer-1}{}}
\newlabel{more-algebra-lemma-finite-cohomology@cref}{{[originallemma][0][]064S}{[1][0][]source}{}{}{}}
\makeatother
\title{The top cohomology of a pseudo-coherent complex with bounded cohomology is finite type and finitely presented in the next degree}
\author{The Stacks Project Authors}
\date{stacks-project:more-algebra.tex:local-snapshot-20261012; GNU FDL 1.2 or later}
\begin{document}
\maketitle
\setcounter{subsection}{1}
\noindent\textit{Editorial note.} All delivered mathematical text is the author's own, copied from the pinned Stacks Project chapter file more-algebra.tex; no statement or proof step was rewritten, reordered or removed.\par\medskip
\noindent\textit{Editorial note.} The chapter file pinned-original-source.tex carries only \textbackslash{}input\{preamble\} on its first line, so no preamble text lives inside the chapter. The retained file sources/stacks-more-algebra/retained-preamble.tex is a mechanical, line-for-line copy of the project's single shared preamble as staged in the snapshot, namely sources/stacks-chow/preamble.tex lines 14-18 (Xy-pic setup) and 154-217 (theorem declarations and mathematical macros). Nothing was added, removed or redefined, no package was added, and the excerpt serves only as preamble, never as a span coordinate.\par\medskip
\noindent\textit{Editorial note.} Both counted spans start at the \textbackslash{}begin line of their environment, so the statement and the author's proof are delivered whole; surrounding prose of the section is deliberately outside the counted spans and is carried as standard prerequisites instead.\par\medskip
\noindent\textit{Editorial note.} The citation fields (source\_id, source\_version, authors, source\_url, license.name, source\_title) are identical to those of the cards already published from this chapter, so that the batch selector treats them as one source instead of deferring them as a same-source clash.\par\medskip
\noindent\textit{Allowed standard prerequisites (editorial).}\begin{itemize}
\item Pseudo-coherence of a complex at a degree, defined through the existence of a bounded complex of finite projective modules mapping to it. definition and surrounding discussion of the same section, source lines 15900-15963. Supplies the complex E and the map alpha that replaces K in the first step of both parts.
\item If the top cohomology of a bounded complex of modules vanishes then the top differential is surjective. used at proof lines 15986-15988 for the case n larger than m. Produces the surjection E\textasciicircum{}\{n-1\} to E\textasciicircum{}n whose splitting removes the top degree.
\item A quotient of a finite projective module by the image of a map from a finite projective module is finitely presented. used at proof lines 15996-16000 for part (2). Converts the cokernel description of H\textasciicircum{}\{m+1\} into finite presentation.
\end{itemize}\medskip
\ifdefined\hypertarget\hypertarget{sourcepointer-1}{}\fi
\noindent\textit{Source reference.} Original lemma 064S: focused-seeds/AGP/sources/stacks-snapshot-20261011/tags.txt, source lines 7876--7876. The lemma's own canonical tag, quoted so that readers can resolve the label inside the delivered text; it points at the counted statement itself and carries no proof content.\par\medskip
\begin{lemma}
\label{lemma-finite-cohomology}
Let $R$ be a ring. Let $K^\bullet$ be a complex of $R$-modules.
Let $m \in \mathbf{Z}$.
\begin{enumerate}
\item If $K^\bullet$ is $m$-pseudo-coherent and $H^i(K^\bullet) = 0$
for $i > m$, then $H^m(K^\bullet)$ is a finite type $R$-module.
\item If $K^\bullet$ is $m$-pseudo-coherent and $H^i(K^\bullet) = 0$
for $i > m + 1$, then $H^{m + 1}(K^\bullet)$ is a finitely presented
$R$-module.
\end{enumerate}
\end{lemma}

\begin{proof}
Proof of (1). Choose a bounded complex $E^\bullet$ of finite projective
$R$-modules and a map $\alpha : E^\bullet \to K^\bullet$ which induces
an isomorphism on cohomology in degrees $> m$ and a surjection in degree $m$.
It is clear that it suffices to prove the result for $E^\bullet$.
Let $n$ be the largest integer such that $E^n \not = 0$.
If $n = m$, then the result is clear.
If $n > m$, then $E^{n - 1} \to E^n$ is surjective as
$H^n(E^\bullet) = 0$. As $E^n$ is finite projective we see that
$E^{n - 1} = E' \oplus E^n$. Hence it suffices to prove the result
for the complex $(E')^\bullet$ which is the same as $E^\bullet$
except has $E'$ in degree $n - 1$ and $0$ in degree $n$.
We win by induction on $n$.

\medskip\noindent
Proof of (2). Choose a bounded complex $E^\bullet$ of finite projective
$R$-modules and a map $\alpha : E^\bullet \to K^\bullet$ which induces
an isomorphism on cohomology in degrees $> m$ and a surjection in degree $m$.
As in the proof of (1) we can reduce to the case that $E^i = 0$ for
$i > m + 1$. Then we see that
$H^{m + 1}(K^\bullet) \cong
H^{m + 1}(E^\bullet) = \Coker(E^m \to E^{m + 1})$
which is of finite presentation.
\end{proof}

\end{document}
